\scanpage{6}
\mspage{1}
\margin{\deleted{\unsure{II}}}
\deleted{
1. \uline{Tores maximaux.}\deleted{\unsure{et centralisateurs}}

\deleted{A partir de maintenant, nous ferons usage des résultats du
Séminaire Chevalley 56/57 sur les Groupes Algébriques. Suivant l'usage
en théorie des groupes algébriques, nous référons à ce Séminaire par le
sigle BIBLE.}

Soit d'abord G un groupe algébrique sur un corps \uline{algébriquement
clos} k. On appelle \uline{tore maximal} de G un \deleted{tore} sous-groupe algébrique
T de G qui est un \uline{tore} (ce qui signifie ici, k étant algébriquement
clos, qu'il est isomorphe à un groupe de la forme $G_m^r$), et qui est
maximal pour cette propriété. Noter que, k étant parfait, $G_{\mathrm{réd}}$ est
un sous-groupe de G, lisse sur k, c'est donc essentiellement un groupe
algébrique au sens de BIBLE. D'ailleurs, tout sous-groupe réduit de G
(tel un tore) est automatiquement un sous-groupe de $\unsure{R|G}_{\mathrm{réd}}$, par suite les
tores maximaux de G coïncident avec les tores maximaux de $G_{\mathrm{réd}}$. Lorsque
G est affine, donc $G_{\mathrm{réd}}$ affine, un théorème fondamental de Borel
nous dit que deux tores maximaux de G sont conjugués \deleted{\illegible{1}} par un
élément de $G(k)$ \inserted{$=G_{\mathrm{réd}}(k)$} (BIBLE 6, th.4 c ), en particulier ils ont même
dimension. \inserted{Nous appellerons} \deleted{L}\inserted{l}a dimension commune des tores maximaux de G \deleted{est alors appelée}
le \uline{rang réductif} de G. Noter d'ailleurs que la restriction G affine
est inutile, comme il résulte d'un théorème \deleted{\unsure{bien}} connu de Chevalley,
affirmant que tout groupe algébrique connexe lisse sur k est une
extension d'une variété abélienne par un groupe affine\deleted{\unsure{un}};
\deleted{nous n'aurons
pas besoin de cette généralisation.}\inserted{cf N°6 $\unsure{\xi|\S}$.}

Dans \inserted{les N°s 1 à 4,} nous nous bornerons \deleted{aux groupes algébriques}
\inserted{le plus souvent aux} préschémas
en groupes affines sur la base.
\deleted{G sur k étant comme ci-dessus,}
\inserted{Soit G un groupe algébrique lisse sur k,} et T
un tore maximal de G, le centralisateur \inserted{$C(T)$} de T dans G sera appelé le
}
\margin{\figure{\begin{tikzpicture}[baseline=(t.base)]
\node[draw,circle,inner sep=5pt] (t) {$G_{\mathrm{réd}}(k)$};\end{tikzpicture}}}
\margin{pas à la ligne}
\margin{à la ligne \figure{\begin{tikzpicture}[baseline]
\draw (0,0)--(.4,0)--(.4,.3)--(.9,.3);\end{tikzpicture}}}
\margin{5}
